% TOP_PROBLEM: {"id":30007558,"problem_number":"LOCAL-30007558","title":"EFX fair allocation existence","category_id":21,"category":{"id":21,"name":"economics","display_name":"Economics","description":"Open economic theory statements, published research agendas, and proposed scoped specifications; question provenance and historical priority are qualified per record.","slug":"economics","order_index":21,"created_at":"2026-10-08T00:00:00Z"},"status":"open","external_url":null,"legacy_ids":[],"published":false,"statement":"Does every finite set of indivisible goods with any number of agents and additive nonnegative valuations admit a complete allocation that removes all envy after deletion of any good positively valued by the potentially envious agent? The deletion requirement explicitly excludes goods that agent values at zero, following the original 2016 EFX convention, often called EFX+.","economics":{"original_id":47,"field":"Economic theory, equilibrium and social choice","kind":"theoretical","formalization_basis":"source_statement","source_status":"explicit_open","priority_status":"earliest_found","priority_note":"The2016 proceedings explicitly state this question and are conventionally credited with introducing EFX in this form. Earlier related near-envy notions exist; absolute earliest priority was not independently established.","earliest_verified_reference":{"authors":["Ioannis Caragiannis","David Kurokawa","Hervé Moulin","Ariel D. Procaccia","Nisarg Shah","Junxing Wang"],"title":"The Unreasonable Fairness of Maximum Nash Welfare","year":2016,"url":"https://eprints.gla.ac.uk/123283/1/123283.pdf","doi":"10.1145/2940716.2940726","locator":"§4.2, Definition 4.4 and following paragraph, proceedings p.317","evidence_summary":"Immediately after defining EFX using positively valued goods, the authors explicitly leave universal complete-allocation existence unsettled."},"current_reference":{"authors":["Thanasis Lianeas","Alkmini Sgouritsa","Minas Marios Sotiriou"],"title":"EFX Allocation in (Multi)Hypergraphs","year":2026,"url":"https://ojs.aaai.org/index.php/AAAI/article/view/38759","doi":"10.1609/aaai.v40i20.38759","locator":"Abstract","evidence_summary":"The abstract explicitly reports that existence for arbitrary additive valuations remains open and proves results for restricted hypergraph instances."},"formalization_note":"Faithful formalization of the explicitly stated2016 question. Naming conventions are stated explicitly; the stronger all-goods version is not silently substituted."},"statusQualification":"A published open mathematical statement was located; its scope and historical priority are qualified in the provenance notes. The cited source dates are the basis of the status; this audit is not an exhaustive proof of nonresolution as of the preparation date. Faithful formalization of the explicitly stated2016 question. Naming conventions are stated explicitly; the stronger all-goods version is not silently substituted.","statusReviewedAt":"2026-10-08"}
% FIRST_POSTED: 2026-10-08
% ENTRY_KIND: statement_only
% !TeX program = lualatex
\documentclass[11pt]{article}
\usepackage{fontspec}
\usepackage[a4paper,margin=25mm]{geometry}
\usepackage{amsmath,amssymb,mathtools}
\usepackage{xurl}
\usepackage[hidelinks]{hyperref}
\newcommand{\sref}[2]{\hyperlink{source:#1:#2}{[S#2]}}
\setlength{\emergencystretch}{3em}
\begin{document}
% problemId:30007558
\section*{EFX fair allocation existence}\label{30007558}
\subsection{Definitions and mathematical statement}
For every positive integer \(n\), let \(N=\{1,\ldots,n\}\) be agents and let \(G\) be any finite set of indivisible goods. Agent \(i\) has an additive nonnegative valuation \(v_i\), specified by numbers \(v_i(g)\geq0\) for \(g\in G\), with \(v_i(S)=\sum_{g\in S}v_i(g)\) for every bundle \(S\subseteq G\). A complete allocation is a tuple \(A=(A_1,\ldots,A_n)\) of pairwise disjoint bundles whose union equals \(G\); empty bundles are allowed. Under the original positive-good convention, require
\[
 \forall i,j\in N\ \forall g\in A_j\text{ with }v_i(g)>0:
 \quad v_i(A_i)\geq v_i(A_j\setminus\{g\}).
\]
Does a complete allocation satisfying these inequalities exist for every such instance? Resolve the universal existence question by proving the assertion for arbitrary \(n\), \(G\), and valuations, or by exhibiting a finite counterexample with a proof that every complete allocation violates an inequality. No computational efficiency requirement is imposed on an existence proof. Zero-valued goods require care: the displayed condition tests only goods positively valued by the potentially envious agent. This is frequently called \(\mathrm{EFX}^{+}\); some later authors reserve EFX or \(\mathrm{EFX}^{0}\) for the stronger condition testing all goods, including zero-valued ones. The source’s Definition 4.4 uses the displayed positive-good convention, and its immediately following paragraph states the complete-allocation existence question.

\par\textbf{Formulation and scope.} Faithful formalization of the explicitly stated2016 question. Naming conventions are stated explicitly; the stronger all-goods version is not silently substituted.

\par\textbf{Open-question provenance.} An explicit question or research agenda was verified at the source locator. This does not by itself certify that every later version remains unresolved \sref{30007558}{1}.

\par\textbf{Historical priority.} The2016 proceedings explicitly state this question and are conventionally credited with introducing EFX in this form. Earlier related near-envy notions exist; absolute earliest priority was not independently established.

\subsection{Short English statement}
Does every finite set of indivisible goods with any number of agents and additive nonnegative valuations admit a complete allocation that removes all envy after deletion of any good positively valued by the potentially envious agent? The deletion requirement explicitly excludes goods that agent values at zero, following the original 2016 EFX convention, often called EFX+.

\subsection{Sources}
\begin{itemize}\raggedright
\item[\textnormal{[S1]}]\hypertarget{source:30007558:1}{} Ioannis Caragiannis, David Kurokawa, Hervé Moulin, Ariel D. Procaccia, Nisarg Shah, Junxing Wang. 2016. The Unreasonable Fairness of Maximum Nash Welfare. §4.2, Definition 4.4 and following paragraph, proceedings p.317.
\url{https://eprints.gla.ac.uk/123283/1/123283.pdf}.
DOI: 10.1145/2940716.2940726.
{\small\textit{Earliest verified question/agenda reference; historical priority qualified below. Immediately after defining EFX using positively valued goods, the authors explicitly leave universal complete-allocation existence unsettled.}}

\item[\textnormal{[S2]}]\hypertarget{source:30007558:2}{} Thanasis Lianeas, Alkmini Sgouritsa, Minas Marios Sotiriou. 2026. EFX Allocation in (Multi)Hypergraphs. Abstract.
\url{https://ojs.aaai.org/index.php/AAAI/article/view/38759}.
DOI: 10.1609/aaai.v40i20.38759.
{\small\textit{Supporting or current literature; see evidence qualification. The abstract explicitly reports that existence for arbitrary additive valuations remains open and proves results for restricted hypergraph instances.}}
\end{itemize}
\end{document}
